% ================================================================================
% MANUAL TECNICO — GAMMA
% Dirigido al administrador de la plataforma. Cubre las secciones restringidas.
% El uso diario del gestor se documenta en el manual de usuario, y la puesta en
% marcha en el manual de instalacion.
%
% Compilar con:  latexmk -pdf manual_tecnico.tex
% ================================================================================
% ================================================================================
% PREAMBULO COMPARTIDO — Manuales de GAMMA
%
% Lo usan los tres manuales (usuario, tecnico e instalacion) para que compartan
% tipografia, paleta y recuadros. Se incluye con % ================================================================================
% PREAMBULO COMPARTIDO — Manuales de GAMMA
%
% Lo usan los tres manuales (usuario, tecnico e instalacion) para que compartan
% tipografia, paleta y recuadros. Se incluye con % ================================================================================
% PREAMBULO COMPARTIDO — Manuales de GAMMA
%
% Lo usan los tres manuales (usuario, tecnico e instalacion) para que compartan
% tipografia, paleta y recuadros. Se incluye con \input{../preambulo} desde la
% carpeta de cada manual.
%
% La paleta sale de las variables de frontend/src/style.css: la interfaz es
% enteramente acromatica, y los manuales tampoco introducen color.
% ================================================================================
\documentclass[11pt, letterpaper]{article}

\usepackage[spanish,mexico]{babel}
\usepackage[utf8]{inputenc}
\usepackage[T1]{fontenc}

% Tipografia: Inter, la mas cercana en TeX Live a la Geist que usa la
% aplicacion. Sin serifas en todo el documento, igual que la interfaz.
\usepackage[sfdefault]{inter}
\usepackage[scaled=0.86]{beramono}

\usepackage[top=2.4cm, bottom=2.4cm, left=2.6cm, right=2.6cm]{geometry}
\usepackage{graphicx}
\usepackage{xcolor}
\usepackage{enumitem}
\usepackage{fancyhdr}
\usepackage{titlesec}
\usepackage{tcolorbox}
% `enhanced` y `borderline west` viven en skins; `breakable` en su propia
% libreria. Sin declararlas, tcolorbox ignora esas claves en silencio y los
% recuadros pierden la barra lateral que los distingue.
\tcbuselibrary{skins,breakable}
\usepackage{booktabs}
\usepackage[hidelinks]{hyperref}

% ── Paleta ──────────────────────────────────────────────────────────────────
% Tomada de las variables de src/style.css. La interfaz es enteramente
% acromatica —todos sus colores son oklch(L 0 0)— asi que el manual tampoco
% introduce color: solo grises, con el texto como unico elemento de contraste.
\definecolor{fg}{HTML}{0A0A0A}        % --foreground
\definecolor{fgsoft}{HTML}{171717}    % --primary
\definecolor{muted}{HTML}{737373}     % --muted-foreground
\definecolor{surface}{HTML}{F5F5F5}   % --muted / --secondary
\definecolor{surfacealt}{HTML}{FAFAFA}% --sidebar
\definecolor{line}{HTML}{E5E5E5}      % --border

\titleformat{\section}{\LARGE\bfseries\color{fg}}{\thesection}{0.7em}{}
\titlespacing*{\section}{0pt}{1.6em}{0.7em}
\titleformat{\subsection}{\large\bfseries\color{fgsoft}}{\thesubsection}{0.6em}{}
\titlespacing*{\subsection}{0pt}{1.3em}{0.4em}
\titleformat{\subsubsection}[runin]{\bfseries\color{fgsoft}}{}{0pt}{}[.]

\pagestyle{fancy}
\fancyhf{}
\lhead{\footnotesize\color{muted}Manual de usuario}
\rhead{\footnotesize\color{muted}\thepage}
\renewcommand{\headrulewidth}{0.4pt}
\renewcommand{\headrule}{\color{line}\hrule height 0.4pt}

\setlength{\parindent}{0pt}
\setlength{\parskip}{0.62em}

% Enlaces y referencias en el mismo gris del cuerpo: sin subrayados ni azules.
\hypersetup{colorlinks=false}

% ── Recuadros ───────────────────────────────────────────────────────────────
% Se distinguen por una barra lateral y el peso del titulo, no por color.
\newtcolorbox{nota}[1][Nota]{
  enhanced, breakable, colback=surfacealt, colframe=surfacealt,
  borderline west={2pt}{0pt}{line},
  coltitle=fg, fonttitle=\bfseries\footnotesize\sffamily,
  title=#1, boxrule=0pt, arc=0pt, outer arc=0pt,
  left=10pt, right=10pt, top=7pt, bottom=7pt, toptitle=1pt, bottomtitle=3pt}

\newtcolorbox{aviso}[1][Importante]{
  enhanced, breakable, colback=surface, colframe=surface,
  borderline west={2pt}{0pt}{fg},
  coltitle=fg, fonttitle=\bfseries\footnotesize\sffamily,
  title=#1, boxrule=0pt, arc=0pt, outer arc=0pt,
  left=10pt, right=10pt, top=7pt, bottom=7pt, toptitle=1pt, bottomtitle=3pt}

% ── Captura de pantalla con marcador de posicion ────────────────────────────
% Uso: \captura{nombre-archivo}{Pie de la figura}
% Si manual/capturas/<nombre>.png existe, la inserta. Si no, dibuja el hueco.
\newcommand{\captura}[2]{%
  \begin{figure}[htbp]
  \centering
  \IfFileExists{capturas/#1.png}{%
    % Ancho completo, pero acotado en alto: una captura de ventana completa
    % puede ser muy alta y desbordaria la pagina. keepaspectratio evita que se
    % deforme cuando manda la restriccion de altura.
    {\color{line}\fboxrule=0.6pt\fboxsep=0pt\fbox{%
      \includegraphics[width=0.94\textwidth, height=0.62\textheight, keepaspectratio]{capturas/#1.png}}}%
  }{%
    \fcolorbox{line}{surfacealt}{%
      \begin{minipage}[c][3.1cm][c]{0.93\textwidth}
        \centering
        {\footnotesize\bfseries\color{muted}CAPTURA PENDIENTE}\\[5pt]
        {\ttfamily\footnotesize\color{fgsoft} capturas/#1.png}\\[4pt]
        {\footnotesize\color{muted} #2}
      \end{minipage}}%
  }
  \caption{#2}
  \label{fig:#1}
  \end{figure}%
}

% Pies de figura discretos
\usepackage{caption}
\captionsetup{font={small,color=muted}, labelfont={bf,color=muted}, skip=6pt}

 desde la
% carpeta de cada manual.
%
% La paleta sale de las variables de frontend/src/style.css: la interfaz es
% enteramente acromatica, y los manuales tampoco introducen color.
% ================================================================================
\documentclass[11pt, letterpaper]{article}

\usepackage[spanish,mexico]{babel}
\usepackage[utf8]{inputenc}
\usepackage[T1]{fontenc}

% Tipografia: Inter, la mas cercana en TeX Live a la Geist que usa la
% aplicacion. Sin serifas en todo el documento, igual que la interfaz.
\usepackage[sfdefault]{inter}
\usepackage[scaled=0.86]{beramono}

\usepackage[top=2.4cm, bottom=2.4cm, left=2.6cm, right=2.6cm]{geometry}
\usepackage{graphicx}
\usepackage{xcolor}
\usepackage{enumitem}
\usepackage{fancyhdr}
\usepackage{titlesec}
\usepackage{tcolorbox}
% `enhanced` y `borderline west` viven en skins; `breakable` en su propia
% libreria. Sin declararlas, tcolorbox ignora esas claves en silencio y los
% recuadros pierden la barra lateral que los distingue.
\tcbuselibrary{skins,breakable}
\usepackage{booktabs}
\usepackage[hidelinks]{hyperref}

% ── Paleta ──────────────────────────────────────────────────────────────────
% Tomada de las variables de src/style.css. La interfaz es enteramente
% acromatica —todos sus colores son oklch(L 0 0)— asi que el manual tampoco
% introduce color: solo grises, con el texto como unico elemento de contraste.
\definecolor{fg}{HTML}{0A0A0A}        % --foreground
\definecolor{fgsoft}{HTML}{171717}    % --primary
\definecolor{muted}{HTML}{737373}     % --muted-foreground
\definecolor{surface}{HTML}{F5F5F5}   % --muted / --secondary
\definecolor{surfacealt}{HTML}{FAFAFA}% --sidebar
\definecolor{line}{HTML}{E5E5E5}      % --border

\titleformat{\section}{\LARGE\bfseries\color{fg}}{\thesection}{0.7em}{}
\titlespacing*{\section}{0pt}{1.6em}{0.7em}
\titleformat{\subsection}{\large\bfseries\color{fgsoft}}{\thesubsection}{0.6em}{}
\titlespacing*{\subsection}{0pt}{1.3em}{0.4em}
\titleformat{\subsubsection}[runin]{\bfseries\color{fgsoft}}{}{0pt}{}[.]

\pagestyle{fancy}
\fancyhf{}
\lhead{\footnotesize\color{muted}Manual de usuario}
\rhead{\footnotesize\color{muted}\thepage}
\renewcommand{\headrulewidth}{0.4pt}
\renewcommand{\headrule}{\color{line}\hrule height 0.4pt}

\setlength{\parindent}{0pt}
\setlength{\parskip}{0.62em}

% Enlaces y referencias en el mismo gris del cuerpo: sin subrayados ni azules.
\hypersetup{colorlinks=false}

% ── Recuadros ───────────────────────────────────────────────────────────────
% Se distinguen por una barra lateral y el peso del titulo, no por color.
\newtcolorbox{nota}[1][Nota]{
  enhanced, breakable, colback=surfacealt, colframe=surfacealt,
  borderline west={2pt}{0pt}{line},
  coltitle=fg, fonttitle=\bfseries\footnotesize\sffamily,
  title=#1, boxrule=0pt, arc=0pt, outer arc=0pt,
  left=10pt, right=10pt, top=7pt, bottom=7pt, toptitle=1pt, bottomtitle=3pt}

\newtcolorbox{aviso}[1][Importante]{
  enhanced, breakable, colback=surface, colframe=surface,
  borderline west={2pt}{0pt}{fg},
  coltitle=fg, fonttitle=\bfseries\footnotesize\sffamily,
  title=#1, boxrule=0pt, arc=0pt, outer arc=0pt,
  left=10pt, right=10pt, top=7pt, bottom=7pt, toptitle=1pt, bottomtitle=3pt}

% ── Captura de pantalla con marcador de posicion ────────────────────────────
% Uso: \captura{nombre-archivo}{Pie de la figura}
% Si manual/capturas/<nombre>.png existe, la inserta. Si no, dibuja el hueco.
\newcommand{\captura}[2]{%
  \begin{figure}[htbp]
  \centering
  \IfFileExists{capturas/#1.png}{%
    % Ancho completo, pero acotado en alto: una captura de ventana completa
    % puede ser muy alta y desbordaria la pagina. keepaspectratio evita que se
    % deforme cuando manda la restriccion de altura.
    {\color{line}\fboxrule=0.6pt\fboxsep=0pt\fbox{%
      \includegraphics[width=0.94\textwidth, height=0.62\textheight, keepaspectratio]{capturas/#1.png}}}%
  }{%
    \fcolorbox{line}{surfacealt}{%
      \begin{minipage}[c][3.1cm][c]{0.93\textwidth}
        \centering
        {\footnotesize\bfseries\color{muted}CAPTURA PENDIENTE}\\[5pt]
        {\ttfamily\footnotesize\color{fgsoft} capturas/#1.png}\\[4pt]
        {\footnotesize\color{muted} #2}
      \end{minipage}}%
  }
  \caption{#2}
  \label{fig:#1}
  \end{figure}%
}

% Pies de figura discretos
\usepackage{caption}
\captionsetup{font={small,color=muted}, labelfont={bf,color=muted}, skip=6pt}

 desde la
% carpeta de cada manual.
%
% La paleta sale de las variables de frontend/src/style.css: la interfaz es
% enteramente acromatica, y los manuales tampoco introducen color.
% ================================================================================
\documentclass[11pt, letterpaper]{article}

\usepackage[spanish,mexico]{babel}
\usepackage[utf8]{inputenc}
\usepackage[T1]{fontenc}

% Tipografia: Inter, la mas cercana en TeX Live a la Geist que usa la
% aplicacion. Sin serifas en todo el documento, igual que la interfaz.
\usepackage[sfdefault]{inter}
\usepackage[scaled=0.86]{beramono}

\usepackage[top=2.4cm, bottom=2.4cm, left=2.6cm, right=2.6cm]{geometry}
\usepackage{graphicx}
\usepackage{xcolor}
\usepackage{enumitem}
\usepackage{fancyhdr}
\usepackage{titlesec}
\usepackage{tcolorbox}
% `enhanced` y `borderline west` viven en skins; `breakable` en su propia
% libreria. Sin declararlas, tcolorbox ignora esas claves en silencio y los
% recuadros pierden la barra lateral que los distingue.
\tcbuselibrary{skins,breakable}
\usepackage{booktabs}
\usepackage[hidelinks]{hyperref}

% ── Paleta ──────────────────────────────────────────────────────────────────
% Tomada de las variables de src/style.css. La interfaz es enteramente
% acromatica —todos sus colores son oklch(L 0 0)— asi que el manual tampoco
% introduce color: solo grises, con el texto como unico elemento de contraste.
\definecolor{fg}{HTML}{0A0A0A}        % --foreground
\definecolor{fgsoft}{HTML}{171717}    % --primary
\definecolor{muted}{HTML}{737373}     % --muted-foreground
\definecolor{surface}{HTML}{F5F5F5}   % --muted / --secondary
\definecolor{surfacealt}{HTML}{FAFAFA}% --sidebar
\definecolor{line}{HTML}{E5E5E5}      % --border

\titleformat{\section}{\LARGE\bfseries\color{fg}}{\thesection}{0.7em}{}
\titlespacing*{\section}{0pt}{1.6em}{0.7em}
\titleformat{\subsection}{\large\bfseries\color{fgsoft}}{\thesubsection}{0.6em}{}
\titlespacing*{\subsection}{0pt}{1.3em}{0.4em}
\titleformat{\subsubsection}[runin]{\bfseries\color{fgsoft}}{}{0pt}{}[.]

\pagestyle{fancy}
\fancyhf{}
\lhead{\footnotesize\color{muted}Manual de usuario}
\rhead{\footnotesize\color{muted}\thepage}
\renewcommand{\headrulewidth}{0.4pt}
\renewcommand{\headrule}{\color{line}\hrule height 0.4pt}

\setlength{\parindent}{0pt}
\setlength{\parskip}{0.62em}

% Enlaces y referencias en el mismo gris del cuerpo: sin subrayados ni azules.
\hypersetup{colorlinks=false}

% ── Recuadros ───────────────────────────────────────────────────────────────
% Se distinguen por una barra lateral y el peso del titulo, no por color.
\newtcolorbox{nota}[1][Nota]{
  enhanced, breakable, colback=surfacealt, colframe=surfacealt,
  borderline west={2pt}{0pt}{line},
  coltitle=fg, fonttitle=\bfseries\footnotesize\sffamily,
  title=#1, boxrule=0pt, arc=0pt, outer arc=0pt,
  left=10pt, right=10pt, top=7pt, bottom=7pt, toptitle=1pt, bottomtitle=3pt}

\newtcolorbox{aviso}[1][Importante]{
  enhanced, breakable, colback=surface, colframe=surface,
  borderline west={2pt}{0pt}{fg},
  coltitle=fg, fonttitle=\bfseries\footnotesize\sffamily,
  title=#1, boxrule=0pt, arc=0pt, outer arc=0pt,
  left=10pt, right=10pt, top=7pt, bottom=7pt, toptitle=1pt, bottomtitle=3pt}

% ── Captura de pantalla con marcador de posicion ────────────────────────────
% Uso: \captura{nombre-archivo}{Pie de la figura}
% Si manual/capturas/<nombre>.png existe, la inserta. Si no, dibuja el hueco.
\newcommand{\captura}[2]{%
  \begin{figure}[htbp]
  \centering
  \IfFileExists{capturas/#1.png}{%
    % Ancho completo, pero acotado en alto: una captura de ventana completa
    % puede ser muy alta y desbordaria la pagina. keepaspectratio evita que se
    % deforme cuando manda la restriccion de altura.
    {\color{line}\fboxrule=0.6pt\fboxsep=0pt\fbox{%
      \includegraphics[width=0.94\textwidth, height=0.62\textheight, keepaspectratio]{capturas/#1.png}}}%
  }{%
    \fcolorbox{line}{surfacealt}{%
      \begin{minipage}[c][3.1cm][c]{0.93\textwidth}
        \centering
        {\footnotesize\bfseries\color{muted}CAPTURA PENDIENTE}\\[5pt]
        {\ttfamily\footnotesize\color{fgsoft} capturas/#1.png}\\[4pt]
        {\footnotesize\color{muted} #2}
      \end{minipage}}%
  }
  \caption{#2}
  \label{fig:#1}
  \end{figure}%
}

% Pies de figura discretos
\usepackage{caption}
\captionsetup{font={small,color=muted}, labelfont={bf,color=muted}, skip=6pt}



\begin{document}

% ── Portada ─────────────────────────────────────────────────────────────────
\begin{titlepage}
\thispagestyle{empty}
\vspace*{5.5cm}
{\fontsize{46}{50}\selectfont\bfseries\color{fg} GAMMA\par}
\vspace{0.5em}
{\large\color{muted} Gobierno Automatizado del Maestro de Materiales\par}

\vspace{2.6cm}
{\color{line}\hrule height 1pt}
\vspace{1.2em}

{\huge\bfseries\color{fg} Manual técnico\par}
\vspace{0.5em}
{\large\color{muted} Guía del administrador de la plataforma\par}

\vfill
{\color{line}\hrule height 0.4pt}
\vspace{1em}
{\footnotesize\color{muted}
Unidad de Energía \quad·\quad \today\par
\vspace{0.4em}
Cubre las secciones restringidas a cuentas administrativas.
El uso diario se documenta en el manual de usuario, y la puesta en marcha en el
manual de instalación.\par}
\vspace{1cm}
\end{titlepage}

\tableofcontents
\newpage

% ============================================================================
\section{Alcance de este manual}

Este documento está dirigido a quien administra GAMMA: la persona que gestiona
las cuentas, vigila la calidad del catálogo y decide cuándo reentrenar el
modelo. Supone que ya conoce el flujo de creación de materiales; si no es así,
empiece por el manual de usuario.

\begin{center}
\begin{tabular}{@{}lp{9.2cm}@{}}
\toprule
\textbf{Manual} & \textbf{Qué cubre} \\
\midrule
Usuario      & El alta de materiales, la exportación a SAP y la ayuda en pantalla. Lo que hace un gestor a diario. \\
Técnico      & Este documento: las cinco secciones restringidas a administradores. \\
Instalación  & Puesta en marcha, actualización y operación de los contenedores. \\
\bottomrule
\end{tabular}
\end{center}

\subsection{Las cinco secciones restringidas}

Un administrador ve en la barra superior cinco secciones que un gestor no:

\begin{center}
\begin{tabular}{@{}llp{7.4cm}@{}}
\toprule
\textbf{Sección} & \textbf{Ruta} & \textbf{Para qué sirve} \\
\midrule
Arquitectura & \texttt{/arquitectura} & Cómo está construida la plataforma \\
Referencia   & \texttt{/referencia}   & Documentación del sistema y de su base de datos \\
Lab          & \texttt{/lab}          & Cuadernos de análisis del proyecto \\
Modelo       & \texttt{/modelo}       & Reentrenamiento y versiones del clasificador \\
Usuarios     & \texttt{/usuarios}     & Alta, baja y modificación de cuentas \\
\bottomrule
\end{tabular}
\end{center}

\newpage
% ============================================================================
\section{Roles y control de acceso}

GAMMA distingue dos roles: \textbf{gestor} y \textbf{administrador}. No hay
niveles intermedios.

\subsection{Qué cambia según el rol}

\begin{center}
\begin{tabular}{@{}lcc@{}}
\toprule
 & \textbf{Gestor} & \textbf{Administrador} \\
\midrule
Chat, Datos y Ayuda                & Sí & Sí \\
Arquitectura, Referencia y Lab     & No & Sí \\
Modelo y Usuarios                  & No & Sí \\
Su actividad cuenta en los indicadores & Sí & No \\
\bottomrule
\end{tabular}
\end{center}

\begin{aviso}[La actividad de un administrador no se contabiliza]
Las cuentas administrativas se excluyen de todas las vistas de indicadores. Fue
una decisión deliberada: el tráfico de pruebas y demostraciones contaminaría las
mediciones de tiempo y de calidad del equipo.

La consecuencia práctica es que \textbf{un administrador no debe usar su cuenta
para dar de alta materiales de la operación real}: esas solicitudes no
aparecerían en ningún reporte. Si necesita crear materiales, hágalo con una
cuenta de gestor.
\end{aviso}

\subsection{Cómo se aplica la restricción}

El control ocurre en tres capas, y solo la última es la que realmente protege:

\begin{enumerate}[leftmargin=*, itemsep=3pt]
  \item \textbf{La barra superior} oculta los enlaces que el rol no permite. Es
        comodidad visual, nada más.
  \item \textbf{La navegación} revalida el rol contra el servidor antes de
        abrir una sección restringida, de modo que escribir la dirección a mano
        no sirve.
  \item \textbf{El servidor} verifica el rol en cada llamada. Aunque alguien
        burlara las dos capas anteriores, no obtendría ningún dato.
\end{enumerate}

\newpage
% ============================================================================
\section{Usuarios}

La sección \textbf{Usuarios} se abre desde el menú de la cuenta, en el extremo
derecho de la barra: haga clic en su nombre y elija \textbf{Administrar}.

\captura{01-menu-administrar}{Menú de la cuenta con la opción Administrar}

\captura{02-usuarios-listado}{Listado de cuentas}

El listado muestra nombre, correo, rol, estado y fecha de alta de cada cuenta
vigente.

\subsection{Crear una cuenta}

Presione \textbf{Nueva cuenta} y complete el formulario. La contraseña debe
tener al menos ocho caracteres. Los dos interruptores del pie definen si la
cuenta es administrativa y si está habilitada.

\captura{03-usuarios-nueva}{Formulario de alta}

\subsection{Editar una cuenta}

El botón \textbf{Editar} de cada fila abre el mismo formulario con los datos
cargados. Puede cambiar nombre, correo, rol y estado.

\begin{nota}[La contraseña se deja en blanco]
Al editar, el campo de contraseña aparece vacío. Déjelo así para conservar la
actual; escriba una nueva solo si quiere cambiarla.
\end{nota}

\subsection{Deshabilitar frente a dar de baja}

Son dos cosas distintas y conviene no confundirlas:

\begin{itemize}[leftmargin=*, itemsep=3pt]
  \item \textbf{Deshabilitar} (quitar la marca de \emph{cuenta habilitada})
        suspende el acceso sin dar de baja a la persona. Úselo para ausencias
        temporales.
  \item \textbf{Dar de baja} retira la cuenta del listado. Úselo cuando alguien
        deja el equipo.
\end{itemize}

\subsection{La baja es lógica}

Dar de baja \textbf{no borra nada}. La cuenta se marca como retirada y
desaparece del listado, pero sus solicitudes, conversaciones y decisiones sobre
duplicados conservan la referencia a quien las creó. El historial no pierde
trazabilidad.

Marque \textbf{Mostrar cuentas dadas de baja} para verlas, junto con la fecha en
que se retiraron y un botón para restaurarlas.

\captura{04-usuarios-bajas}{Cuentas dadas de baja}

\begin{nota}[El correo se libera]
Al dar de baja una cuenta, su correo queda disponible para volver a usarse. Si
después crea una cuenta nueva con ese mismo correo, la antigua ya no se podrá
restaurar mientras la nueva exista: el sistema se lo advertirá.
\end{nota}

\subsection{Resguardos que no puede saltarse}

El sistema impide cuatro operaciones que dejarían la plataforma sin
administración:

\begin{itemize}[leftmargin=*, itemsep=3pt]
  \item Retirarse a usted mismo el rol de administrador.
  \item Desactivar su propia cuenta.
  \item Darse de baja a usted mismo.
  \item Dejar el sistema sin ningún administrador activo.
\end{itemize}

En todos los casos aparece un mensaje explicando el motivo y la operación no se
realiza. Si necesita retirarse, promueva primero a otra persona.

\newpage
% ============================================================================
\section{Modelo}

La sección \textbf{Modelo} controla el clasificador que sugiere la categoría de
cada material. Se llega desde la barra superior.

\captura{05-modelo-pantalla}{Pantalla de modelo}

\subsection{Cuándo reentrenar}

Casi nunca. El reentrenamiento tiene sentido después de una \textbf{carga
importante al maestro}: materiales nuevos, clases nuevas o una corrección
masiva del catálogo. No es una tarea programada ni periódica.

\begin{aviso}[Es un proceso costoso]
Tarda alrededor de diez minutos y mantiene en memoria el modelo vigente y el
nuevo a la vez. Durante ese lapso el servicio de predicción puede responder más
lento. Ejecútelo fuera del horario de mayor uso.
\end{aviso}

\subsection{Cómo se ejecuta}

\begin{enumerate}[leftmargin=*, itemsep=3pt]
  \item Presione \textbf{Iniciar reentrenamiento} y confirme.
  \item La pantalla muestra la etapa en curso y se actualiza sola cada pocos
        segundos. Puede cerrarla: el proceso sigue en el servidor.
  \item Al terminar aparece la comparación de métricas contra el modelo
        anterior.
\end{enumerate}

\captura{06-modelo-comparacion}{Comparación de métricas al terminar}

\subsection{Cómo leer la comparación}

Se comparan cuatro medidas. La diferencia se expresa en puntos porcentuales:
verde si mejoró, rojo si empeoró.

\begin{center}
\begin{tabular}{@{}lp{9.2cm}@{}}
\toprule
\textbf{Medida} & \textbf{Qué significa} \\
\midrule
Exactitud     & Proporción de aciertos en la primera sugerencia \\
F1 macro      & Desempeño promediado dando el mismo peso a cada clase, sin importar cuántos materiales tenga \\
F1 ponderado  & Lo mismo, pero pesando cada clase por su tamaño \\
Top-3         & Proporción de casos donde la clase correcta aparece entre las tres sugerencias \\
\bottomrule
\end{tabular}
\end{center}

\begin{nota}[Fíjese en el F1 macro]
Si la exactitud sube pero el F1 macro baja, el modelo mejoró en las clases
grandes a costa de las pequeñas. Como buena parte del catálogo son clases con
pocos materiales, esa combinación no siempre es una mejora real.
\end{nota}

\subsection{El modelo nuevo entra de inmediato}

Al terminar, el modelo reentrenado queda en servicio sin necesidad de reiniciar
nada. Por eso conviene revisar la comparación en cuanto termina: si algo
empeoró de forma relevante, revierta.

\subsection{Versiones y reversión}

Cada reentrenamiento archiva el modelo anterior con una marca de tiempo antes de
publicar el nuevo. El listado de versiones muestra las disponibles con sus
métricas y su tamaño.

\captura{07-modelo-versiones}{Historial de versiones}

Para revertir, presione \textbf{Revertir} en la versión que quiera restaurar y
confirme. El modelo vigente se archiva también, de modo que una reversión
puede deshacerse.

\begin{nota}[Se conservan las tres últimas]
Los artefactos ocupan espacio considerable en disco, así que solo se guardan las
tres versiones más recientes. Una versión listada sin archivo disponible aparece
señalada y ya no se puede restaurar.
\end{nota}

\subsection{Con qué datos se reentrena}

El conjunto se arma con el maestro de materiales vigente más las solicitudes que
un gestor \textbf{corrigió de forma explícita}.

Quedan fuera, a propósito, las solicitudes que el modelo resolvió por su cuenta
y el gestor aceptó sin cambios: reentrenar con las propias sugerencias aceptadas
refuerza los sesgos del modelo en lugar de corregirlos. También se descartan las
clases con muy pocos materiales, insuficientes para entrenar y evaluar.

\begin{aviso}[Una clase nueva necesita ejemplos]
Un material cuya clase acaba de crearse no será sugerido por el modelo hasta que
esa clase tenga al menos tres materiales en el maestro y se haya reentrenado.
\end{aviso}

\newpage
% ============================================================================
\section{Arquitectura}

La sección \textbf{Arquitectura} es documentación de consulta: describe cómo
está construida la plataforma. No tiene controles ni acciones.

\captura{08-arquitectura}{Sección de arquitectura}

Contiene el diagrama de los componentes y su relación, la tabla del stack
tecnológico con la justificación de cada elección, y la descripción de la
arquitectura de datos por capas.

Es útil para explicar el sistema a alguien de fuera, o para ubicarse antes de
tocar algo.

\newpage
% ============================================================================
\section{Referencia}

La sección \textbf{Referencia} es la documentación técnica detallada,
organizada en pestañas.

\captura{09-referencia}{Sección de referencia}

\begin{center}
\begin{tabular}{@{}lp{9.2cm}@{}}
\toprule
\textbf{Pestaña} & \textbf{Contenido} \\
\midrule
API             & Documentación interactiva de todas las rutas. Permite probarlas desde el navegador \\
Esquema         & Diagramas de la base de datos por capa, con sus tablas y relaciones \\
Decisiones      & Las decisiones de diseño del proyecto y su justificación \\
Modelo          & Cómo funciona el clasificador: representación del texto, algoritmo y métricas \\
Reentrenamiento & Detalle técnico del mecanismo de versionado y las rutas que lo exponen \\
\bottomrule
\end{tabular}
\end{center}

\begin{nota}[La pestaña API es interactiva]
No es solo documentación: permite ejecutar llamadas reales contra el sistema.
Tenga presente que lo que ejecute ahí \textbf{afecta los datos de verdad}.
\end{nota}

\newpage
% ============================================================================
\section{Lab}

La sección \textbf{Lab} reúne los cuadernos de análisis del proyecto. Son de
consulta: documentan cómo se llegó a los resultados, no forman parte de la
operación.

\captura{10-lab}{Sección de laboratorio}

Se agrupan en dos bloques:

\begin{itemize}[leftmargin=*, itemsep=3pt]
  \item \textbf{Análisis exploratorio}, con las vistas consolidadas del maestro
        y un cuaderno por cada tipo de material, con su volumen y su tasa de
        duplicados.
  \item \textbf{Modelado y analítica de negocio}, con la comparación de
        algoritmos y el análisis de ahorro y retorno de la inversión.
\end{itemize}

El botón \textbf{Jupyter} de la esquina abre el entorno donde los cuadernos se
pueden ejecutar y modificar.

\newpage
% ============================================================================
\section{Operación recomendada}

\subsection{Al incorporar a alguien al equipo}

Cree la cuenta desde \textbf{Usuarios} sin marcar el interruptor de
administrador. Entréguele el manual de usuario. Marque administrador solo si esa
persona va a gestionar cuentas o el modelo.

\subsection{Al actualizar el maestro}

Cargue los archivos desde \textbf{Datos}, pestaña Importar, respetando el orden:
UNSPSC, luego clases y al final materiales. Verifique que los conteos hayan
subido. Reentrene solo si la carga fue significativa.

\subsection{Al reentrenar}

Hágalo fuera del horario de mayor uso, revise la comparación de métricas en
cuanto termine, y revierta si alguna empeoró de forma relevante.

\subsection{Revisión periódica}

Vale la pena mirar de vez en cuando los indicadores de la sección Datos: si la
tasa de corrección del modelo sube de forma sostenida, es señal de que el
catálogo cambió lo suficiente como para justificar un reentrenamiento.

\vspace{1cm}

\begin{nota}[Puesta en marcha y actualización]
Todo lo relativo a instalar, actualizar o resolver problemas de los contenedores
está en el manual de instalación.
\end{nota}

\end{document}
